\documentclass[11pt,a4paper]{article}
\usepackage[T1]{fontenc}
\usepackage[utf8]{inputenc}
\usepackage[main=italian,provide=*]{babel}
\usepackage{amsmath,amssymb}
\usepackage{geometry}
\geometry{margin=2.3cm,top=2.5cm,bottom=2.7cm}
\usepackage{booktabs}
\usepackage[table,dvipsnames]{xcolor}
\usepackage{tcolorbox}
\tcbuselibrary{skins,breakable}
\usepackage{titlesec}
\usepackage{enumitem}
\usepackage{needspace}
\usepackage{fancyhdr}
\usepackage{graphicx}
\usepackage{hyperref}
\hypersetup{colorlinks=true,linkcolor=Sepia,citecolor=Sepia,urlcolor=Sepia}
\newcommand{\Tr}{\operatorname{Tr}}
\newcommand{\ket}[1]{\lvert #1 \rangle}
\newcommand{\bra}[1]{\langle #1 \rvert}

\definecolor{inkblue}{HTML}{1B3A5C}
\definecolor{lightblue}{HTML}{EAF1F8}
\definecolor{accent}{HTML}{B0562A}
\definecolor{softgray}{HTML}{6B6B6B}
\definecolor{okgreen}{HTML}{2E6B3E}
\definecolor{okgreenbg}{HTML}{EAF5EC}

\titleformat{\section}
  {\normalfont\Large\bfseries\color{inkblue}}
  {\thesection}{0.6em}{}[{\color{inkblue!40}\titlerule}]
\titlespacing{\section}{0pt}{0.9em}{0.4em}
\titleformat{\subsection}
  {\normalfont\large\bfseries\color{accent}}
  {}{0em}{}
\titlespacing{\subsection}{0pt}{0.6em}{0.2em}

\pagestyle{fancy}
\fancyhf{}
\renewcommand{\headrulewidth}{0.4pt}
\fancyhead[L]{\small\color{softgray}VQE noise-aware sul dimero --- resoconto definitivo}
\fancyhead[R]{\small\color{softgray}\thepage}
\fancyfoot[C]{\small\color{softgray}Tesi triennale, Samuele --- Relatore: Prof.\ Alessandro Chiesa}

\newtcolorbox{keyeq}[1][]{
  colback=lightblue, colframe=inkblue!70, boxrule=0.6pt,
  arc=2pt, left=8pt, right=8pt, top=3pt, bottom=3pt, #1
}
\newtcolorbox{okbox}[1]{
  colback=okgreenbg, colframe=okgreen!60, boxrule=0.6pt,
  arc=2pt, left=7pt, right=7pt, top=2pt, bottom=2pt,
  fonttitle=\bfseries\color{okgreen}, title={#1}
}

\title{\vspace{-1.5em}\color{inkblue}VQE noise-aware sul dimero\\[0.2em]
\Large\normalfont\color{softgray}Resoconto definitivo: obiettivo, metodo, risultati, generalità}
\author{Note di lavoro --- Tesi triennale, Samuele\\ Relatore: Prof.\ Alessandro Chiesa}
\date{}

\begin{document}
\sloppy
\maketitle
\thispagestyle{fancy}
\vspace{-2.0em}

File di codice: \texttt{vqe\_noise\_aware\_dimero.py},
\texttt{validate\_vqe\_noise\_aware\_dimero.py},
\texttt{verifica\_passo5\_noise\_aware.py}. Punto di lavoro principale:
$b/J=0.35$, $D/J=0.80$, $J=1$ (``test 2''); un secondo punto,
$b/J=-0.18$, $D/J=1$, verifica la generalità (§\ref{sec:secondo-punto}).
Modello di rumore: \texttt{noise\_model\_dimero.py} (calibrazione reale
\texttt{ibm\_torino}, arXiv:2504.15187).

\tableofcontents

\section{Obiettivo}

La pipeline di rumore (Parte 2) misura, nel VQE con termine DM sotto
rumore (Passo 2), quanto degrada un
risultato VQE ideale quando lo si esegue su un simulatore rumoroso: i
parametri $\theta^*_\text{ideale}$ sono ottimizzati \textbf{senza}
rumore, poi il circuito con quella ricetta viene valutato \textbf{con}
rumore. Questo documento affronta la domanda complementare: se il rumore
è acceso già durante l'ottimizzazione --- come avverrebbe su un hardware
reale, dove non esiste una versione ideale a cui accedere per calibrare
l'ottimizzatore --- l'ottimizzatore converge a parametri diversi? E se
sì, questo cambia in modo misurabile i risultati a valle (Trotter,
correlatori dinamici, scan sui parametri di rumore)?

È un'estensione \textbf{facoltativa} della pipeline (già chiusa in tutti
e cinque i passi), non uno step mancante. Lo scope copre: la
rioptimizzazione stessa (Passo 2), il suo effetto su $N^*$ della
fedeltà di Trotter (Passo 3) e su $N^*$ del correlatore dinamico (Passo
4), la sua generalità rispetto al punto di lavoro fisico e ai parametri
di rumore (Passo 5, griglia completa).

\section{Argomento fisico e previsione}
\label{sec:argomento}

Il canale di rumore è depolarizzante, isotropo: riduce uniformemente la
sfera di Bloch in ogni direzione. Per un canale di questo tipo, l'azione
su uno stato puro $\ket{\psi(\theta)}$ preparato dall'ansatz si scrive
schematicamente
\begin{keyeq}
\centering
$\rho(\theta) = (1-\lambda)\,\ket{\psi(\theta)}\bra{\psi(\theta)} + \lambda\,\dfrac{\mathbf{1}}{d}$
\end{keyeq}
con $\lambda>0$ legato a $\varepsilon_{1q},\varepsilon_{2q}$ tramite le
relazioni derivate in \texttt{teoria\_modello\_rumore.tex}, $d=4$ per 2
qubit. L'energia rumorosa diventa
\begin{keyeq}
\centering
$E_\text{rumoroso}(\theta) = \Tr[\rho(\theta) H] = (1-\lambda)\,E_\text{ideale}(\theta) + \lambda\,\dfrac{\Tr[H]}{d}$
\end{keyeq}
--- una trasformazione \textbf{affine} di $E_\text{ideale}(\theta)$: il
secondo termine non dipende da $\theta$, il primo è $E_\text{ideale}(\theta)$
moltiplicato per una costante positiva $1-\lambda$. Una trasformazione
affine con coefficiente positivo non cambia la posizione del minimo:
\textbf{lo stesso $\theta^*$ che minimizza $E_\text{ideale}$ minimizza
anche $E_\text{rumoroso}$}. A questo si aggiunge una seconda osservazione,
verificata indipendentemente (§\ref{sec:verifica-cnot}): il conteggio CNOT
del circuito transpilato non dipende dai valori numerici dei parametri,
quindi non c'è nemmeno una via indiretta (via costo in gate) per cui
convenga spostarsi verso un $\theta$ diverso da quello ideale.

\begin{keyeq}[title={Previsione dichiarata prima della verifica}]
\centering
$\theta^*_\text{rumore} \approx \theta^*_\text{ideale}$: nessun vantaggio
misurabile a rioptimizzare sotto rumore rispetto al semplice riuso dei
parametri ideali.
\end{keyeq}

Questa è una previsione, non un fatto acquisito: il resto del documento
la mette alla prova, su più metriche e più punti di lavoro.

\subsection{Due metriche a valle, strutturalmente diverse}
\label{sec:due-metriche}

La previsione riguarda l'\emph{energia} $E(\theta)$. Ma l'effetto di
$\theta^*_\text{rumore}\approx\theta^*_\text{ideale}$ va controllato anche
sulle quantità che dipendono dalla preparazione \emph{a valle} nella
pipeline --- ed è qui che serve una distinzione. La quantum simulation
(Trotter) sotto rumore (Passo 3) usa un osservabile additivo, $F(N)=\bra{\psi_\text{target}}\rho_N\ket{\psi_\text{target}}$;
i correlatori dinamici sotto rumore (Passo 4) usano invece $N^*=\arg\max_N|C(N)|$, il
\textbf{modulo} di un numero complesso costruito da una misura di
Hadamard test. Sono oggetti matematici diversi: un conto valido per una
somma additiva non si trasporta automaticamente a un modulo (si veda
\texttt{teoria\_readout\_asimmetrico.tex} per un caso concreto in cui
questa distinzione cambia una conclusione). Per questo entrambe le
metriche vengono verificate esplicitamente in questo documento, non solo
una delle due per analogia.

\section{Implementazione}

\subsection{Ansatz e ottimizzazione}

L'ansatz è il PMA-2q$\cdot$3 di Parte 1 (blocco RBS-conservante il
numero di eccitazioni, più due rotazioni $R_y$ indipendenti, 3
parametri). L'energia rumorosa, per un dato $\theta$, si calcola
simulando il circuito su \texttt{AerSimulator(method="density\_matrix")}
con il \texttt{NoiseModel} agganciato:
\begin{keyeq}
\centering
$E_\text{rumoroso}(\theta) = \Tr[\rho(\theta) H], \qquad \rho(\theta) = \mathcal{N}\!\big(\ket{\psi(\theta)}\bra{\psi(\theta)}\big)$
\end{keyeq}
Poiché la simulazione a operatore densità è \textbf{esatta} (nessun
campionamento a shot finiti), $E_\text{rumoroso}(\theta)$ è liscia in
$\theta$: gli stessi ottimizzatori del caso ideale (COBYLA multistart +
polish L-BFGS-B) restano applicabili senza modifiche concettuali.

\subsection{Una scelta tecnica: transpilare dopo, non prima, dell'assegnazione}
\label{sec:trappola}

Il circuito va transpilato \textbf{dopo} aver assegnato i valori numerici
di $\theta$, ad ogni valutazione dell'obiettivo --- non transpilato una
volta sola in forma parametrica e poi solo legato con
\texttt{assign\_parameters} (il principio che invece funziona, e va
usato, per il blocco di Trotter ripetuto $N$ volte, dove i parametri
fisici non cambiano da una valutazione all'altra). La ragione: con
parametri simbolici il transpilatore non può fondere le rotazioni a 1
qubit il cui angolo dipende dal valore numerico di $\theta$, producendo
un circuito strutturalmente più rumoroso (Fig.~\ref{fig:bug}: 19
\texttt{rz} + 14 \texttt{sx} contro 11 \texttt{rz} + 8 \texttt{sx}, a
parità di 2 CNOT). Più costoso per valutazione ($\sim9\,$ms), ma l'unico
approccio che riproduce il conteggio di gate minimo già stabilito nella
pipeline principale.

\begin{figure}[h]
\centering
\includegraphics[width=0.6\linewidth]{figure/fig_bug_conteggio_gate.pdf}
\caption{\textit{Costo in gate del circuito transpilato, a parità di
parametri $\theta$, a seconda dell'ordine fra transpilazione e
assegnazione. Il conteggio CNOT resta invariato (2); i gate a singolo
qubit quasi raddoppiano se si transpila prima dell'assegnazione --- ogni
gate in più è un'ulteriore occasione di errore per il canale
depolarizzante.}}
\label{fig:bug}
\end{figure}

\section{Verifica: invarianza del conteggio CNOT}
\label{sec:verifica-cnot}

Premessa fisica di §\ref{sec:argomento}: il conteggio CNOT del circuito
transpilato (nell'ordine corretto) non deve dipendere dai valori
numerici di $\theta$. Verificato su 5 assegnazioni casuali indipendenti:
CNOT$=2$ in tutte le prove (11 \texttt{rz} + 8 \texttt{sx} + 2
\texttt{cx} in ciascuna) --- nessuna direzione del paesaggio dei
parametri verso cui convenga spostarsi per ``risparmiare'' rumore.

\section{Risultato 1: energia e fedeltà (Passo 2)}

Multistart $R=6$ punti casuali $+\,1$ seme uguale a $\theta^*_\text{ideale}$
(per garantire di non fare mai peggio del semplice riuso; la robustezza
del multistart è discussa in §\ref{sec:robustezza}).

\begin{center}
\renewcommand{\arraystretch}{1.25}
\begin{tabular}{@{}lcc@{}}
\toprule
\textbf{Quantità} & \textbf{Riuso $\theta^*_\text{ideale}$} & \textbf{VQE noise-aware}\\
\midrule
$E$ (rumoroso) & $-3.50164605$ & $-3.50164616$\\
$F$ (rumorosa) & $0.98500601$ & $0.98500602$\\
CNOT & $2$ & $2$\\
\bottomrule
\end{tabular}
\end{center}

\noindent $\Delta E=-1.16\times10^{-7}$, $\Delta F=+1.76\times10^{-8}$ ---
sei ordini di grandezza più piccoli della scala fisica del problema
(perdita di fedeltà $\sim1.5\times10^{-2}$ a rumore acceso). Lo stesso
confronto, ripetuto sull'intero scan di $\varepsilon_{2q}$ (Fig.~\ref{fig:eps2q}),
mostra le due curve sovrapposte su tutto il range esplorato.

\begin{figure}[h]
\centering
\includegraphics[width=0.92\linewidth]{figure/fig_scan_eps2q_confronto.pdf}
\caption{\textit{Energia e fedeltà rumorose in funzione di
$\varepsilon_{2q}$, riuso dei parametri ideali contro rioptimizzazione
noise-aware. Le due curve sono indistinguibili a questa scala su tutto
il range, dal rumore nullo a $4\times$ il valore di riferimento.}}
\label{fig:eps2q}
\end{figure}

\section{Robustezza del multistart}
\label{sec:robustezza}

Con $R=6$ punti casuali e nessun seme, il minimo trovato può restare a
distanza $10^{-4}$--$10^{-2}$ dall'ottimo vero, a seconda del seed
(dispersione reale del paesaggio energetico rumoroso, caratterizzata su
8 seed indipendenti, Fig.~\ref{fig:multistart}). Per questo il risultato
principale include $\theta^*_\text{ideale}$ come uno dei punti di
partenza. Un multistart indipendente più ampio ($R=20$, nessun seme)
converge comunque allo stesso ottimo entro $2.6\times10^{-6}$: il seme
non altera la conclusione fisica, garantisce solo di raggiungerla con
un numero contenuto di restart.

\begin{figure}[h]
\centering
\includegraphics[width=0.72\linewidth]{figure/fig_robustezza_multistart.pdf}
\caption{\textit{Energie finali dei singoli restart COBYLA (rumore
acceso), per tre configurazioni di multistart. La linea tratteggiata è
l'ottimo noto. Con soli 6 restart casuali il minimo trovato può restare
lontano dall'ottimo; con 20 il punto migliore si avvicina entro
$3.6\times10^{-6}$.}}
\label{fig:multistart}
\end{figure}

\section{Risultato 2: \texorpdfstring{$N^*$}{N*} della fedeltà di Trotter (Passo 3)}

\begin{center}
\renewcommand{\arraystretch}{1.2}
\begin{tabular}{@{}lcc@{}}
\toprule
\textbf{Preparazione} & $\boldsymbol{N^*}$ & $\boldsymbol{F(N^*)}$\\
\midrule
$\theta^*_\text{ideale}$ & $8$ & $0.817725$\\
$\theta^*_\text{noise-aware}$ & $8$ & $0.817767$\\
\bottomrule
\end{tabular}
\end{center}

\begin{figure}[h]
\centering
\includegraphics[width=0.92\linewidth]{figure/fig_F_vs_N_confronto.pdf}
\caption{\textit{Sinistra: fedeltà $F(N)$, le due preparazioni a
confronto --- massimo a $N^*=8$ in entrambi i casi. Destra: scostamento,
scala ingrandita ($\sim10^{-4}$): un ordine di grandezza sopra il rumore
numerico, ma sei ordini sotto la scala fisica del problema.}}
\label{fig:FvsN}
\end{figure}

\section{Risultato 3: \texorpdfstring{$N^*$}{N*} del correlatore dinamico (Passo 4)}

Per la distinzione strutturale di §\ref{sec:due-metriche}, questa
metrica è verificata \textbf{direttamente}, non per estensione dal
risultato precedente.

\begin{center}
\renewcommand{\arraystretch}{1.2}
\begin{tabular}{@{}lcc@{}}
\toprule
\textbf{Preparazione} & $\boldsymbol{N^*}$ & $\boldsymbol{|C(N^*)|}$\\
\midrule
$\theta^*_\text{ideale}$ & $5$ & $0.250002$\\
$\theta^*_\text{noise-aware}$ & $5$ & $0.249933$\\
\bottomrule
\end{tabular}
\end{center}

\begin{figure}[h]
\centering
\includegraphics[width=0.92\linewidth]{figure/fig_correlatore_noise_aware_confronto.pdf}
\caption{\textit{Sinistra: modulo del correlatore rumoroso, le due
preparazioni a confronto --- massimo a $N^*=5$ in entrambi i casi.
Destra: scostamento, scala ingrandita ($\sim10^{-5}$) --- un ordine di
grandezza più piccolo persino dello scostamento osservato per la
fedeltà di Trotter.}}
\label{fig:corrNA}
\end{figure}

\section{Perché il risultato generalizza: l'argomento di continuità}
\label{sec:continuita}

Le due metriche condividono una proprietà comune, verificabile
direttamente: $\theta^*_\text{noise-aware}$ e $\theta^*_\text{ideale}$
sono quasi indistinguibili ($\Delta E\sim10^{-7}$). Per qualunque
funzione ragionevolmente liscia dello stato preparato --- una fedeltà,
il modulo di un correlatore, o altro --- questo garantisce, per
continuità, uno scostamento a valle altrettanto minuscolo,
indipendentemente dalla struttura specifica della metrica (additiva o
basata su modulo). Non è quindi la struttura di $N^*$ a garantire
l'invarianza: è la piccolezza di $\Delta\theta$ a renderla ininfluente su
qualunque metrica ragionevole a valle.

Questo argomento è stato reso quantitativo e verificato (non solo
asserito): lo scostamento osservato in entrambe le metriche è riprodotto
entro lo $0.1\%$ da una previsione lineare al prim'ordine,
$\Delta f\approx\nabla f(\theta^*_\text{ideale})\cdot\Delta\theta$ ---
derivazione e verifica complete in una nota di lavoro interna
(\texttt{nota\_continuita\_preparazione.tex}, non riportata per esteso
qui per la sua natura metodologica più che di risultato fisico).

\section{Generalità: un secondo punto di lavoro}
\label{sec:secondo-punto}

L'argomento fisico di §\ref{sec:argomento} è generale --- vale per
qualunque $b/J,D/J$ e qualunque intensità di rumore isotropo --- ma la
\emph{dimensione} di $\Delta\theta$, e quindi se è abbastanza piccola da
non spostare un $N^*$ specifico, è un fatto numerico. Per non lasciarlo
verificato in un solo punto, l'intera catena (VQE ideale, VQE
noise-aware, $N^*$ Trotter, $N^*$ correlatore) è stata ripetuta a un
secondo punto di lavoro: $b/J=-0.18$, $D/J=1$ --- un termine DM cinque
volte più forte e segno di $b$ opposto rispetto a ``test 2'', già
validato con lo stesso ansatz in altre parti del progetto.

\begin{center}
\footnotesize
\renewcommand{\arraystretch}{1.25}
\begin{tabular}{@{}lcc@{}}
\toprule
\textbf{Quantità} & \textbf{test 2} & \textbf{secondo punto}\\
 & ($b/J{=}0.35,D/J{=}0.80$) & ($b/J{=}{-}0.18,D/J{=}1$)\\
\midrule
$\Delta E$ (noise-aware $-$ riuso ideale) & $-1.16\times10^{-7}$ & $-2.56\times10^{-8}$\\
$\Delta F$ (noise-aware $-$ riuso ideale) & $+1.76\times10^{-8}$ & $+3.96\times10^{-9}$\\
$N^*$ Trotter (ideale $\to$ noise-aware) & $8\to8$ & $8\to8$\\
$N^*$ correlatore (ideale $\to$ noise-aware) & $5\to5$ & $3\to3$\\
\bottomrule
\end{tabular}
\end{center}

\begin{figure}[h]
\centering
\includegraphics[width=0.92\linewidth]{figure/fig_secondo_punto_confronto.pdf}
\caption{\textit{Secondo punto di lavoro. Sinistra: fedeltà di Trotter,
$N^*=8$ --- stesso valore del punto ``test 2''. Destra: modulo del
correlatore, $N^*=3$ --- diverso dal punto ``test 2'' ($N^*$ dipende
dalla fisica del punto, non è universale), ma comunque invariato fra le
due preparazioni.}}
\label{fig:secondopunto}
\end{figure}

\begin{okbox}{Cosa è generale e cosa resta specifico di un punto}
$N^*$ \textbf{cambia} da un punto di lavoro all'altro: è una proprietà
della fisica del punto, non della preparazione VQE. Quello che
\textbf{non} cambia, in entrambi i punti testati, è che $N^*$ resta lo
stesso confrontando le due preparazioni --- coerente con l'argomento di
continuità (§\ref{sec:continuita}): $\Delta\theta$ resta piccolo anche a
questo secondo punto, pur essendo un regime fisicamente diverso. Due
punti non dimostrano una legge generale per ogni $b/J,D/J$, ma escludono
che il risultato sia un caso isolato.
\end{okbox}

\section{Generalità: tutta la griglia dello scan sui parametri di rumore}
\label{sec:passo5}

Lo scan sui parametri di rumore (Passo 5) esplora una griglia
$(\varepsilon_{1q},\varepsilon_{2q},p_\text{readout})$. Le due metriche a
valle sono state ricalcolate con \textbf{entrambe} le preparazioni su
\textbf{tutta} la griglia già usata nello scan originale --- non solo al
punto di riferimento --- per un totale di 25 combinazioni.

\begin{figure}[h]
\centering
\includegraphics[width=0.95\linewidth]{figure/fig_passo5_griglia_completa.pdf}
\caption{\textit{$N^*$ in funzione di $\varepsilon_{2q}$, riuso ideale
contro noise-aware, su tutta la griglia dello scan sui parametri di
rumore. Sinistra: fedeltà
di Trotter (con la nota transizione a $N^*=1$ per rumore forte, già
caratterizzata nello scan originale). Destra: modulo del correlatore.
Le due curve sono sovrapposte in ogni punto.}}
\label{fig:passo5}
\end{figure}

\begin{okbox}{25/25}
$N^*$ resta invariato fra le due preparazioni in \textbf{tutti} i 25
punti testati (7 valori di $\varepsilon_{2q}$ e 6 di $\varepsilon_{1q}$
per la fedeltà di Trotter; 5 valori di $\varepsilon_{2q}$, 4 di
$\varepsilon_{1q}$ e 3 di $p_\text{readout}$ per il correlatore) ---
verificato punto per punto, non solo al riferimento.
\end{okbox}

\section{Considerazioni}

\begin{itemize}[leftmargin=1.3em,itemsep=0.35em]
\item \textbf{La previsione è confermata, non solo plausibile.}
L'argomento della trasformazione affine non garantiva a priori uno
scostamento numerico trascurabile su \emph{ogni} metrica a valle
(§\ref{sec:due-metriche}): il modulo di un correlatore e una fedeltà
additiva sono oggetti diversi, e comporre una trasformazione con un
modulo non eredita automaticamente le stesse proprietà di invarianza
(lo si vede in modo netto nel lavoro sul readout asimmetrico,
\texttt{teoria\_readout\_asimmetrico.tex}). Per questo entrambe le
metriche sono state verificate esplicitamente, non una per analogia
dall'altra.
\item \textbf{L'argomento di continuità è più solido di un argomento di
struttura}, e più generale: si applica a qualunque osservabile liscio
dello stato preparato, additivo o no, purché il gradiente resti moderato
(verificato quantitativamente, §\ref{sec:continuita}) --- non richiede
che l'osservabile abbia una decomposizione additiva nota a priori.
\item \textbf{Il dominio di validità è dichiarato esplicitamente}: due
punti di lavoro, un insieme di parametri di rumore esplorato per intero
al punto di riferimento del canale. Non è una dimostrazione per ogni
$b/J,D/J$ o ogni intensità di rumore, ma un controllo sistematico, non
un singolo punto isolato.
\end{itemize}

\section{Conclusioni}

Rioptimizzare l'ansatz PMA-2q$\cdot$3 sotto rumore, invece di riusare i
parametri ideali, non cambia in modo misurabile né l'energia/fedeltà del
Passo 2 ($\Delta E\sim10^{-7}$--$10^{-8}$ a seconda del punto di lavoro),
né la posizione di $N^*$ della fedeltà di Trotter (Passo 3, $N^*=8$
invariato) o del correlatore dinamico (Passo 4, $N^*=5$ o $N^*=3$ a
seconda del punto di lavoro, comunque invariato). Il risultato è
verificato --- non solo argomentato per analogia --- su due metriche
strutturalmente diverse, due punti di lavoro fisici, e sull'intera
griglia di parametri di rumore del Passo 5 (25/25 punti). L'argomento
che lo spiega, la continuità dell'osservabile in funzione di un
$\Delta\theta$ minuscolo, è stato reso quantitativo e verificato entro lo
$0.1\%$.

\begin{okbox}{Sintesi in una riga}
Rioptimizzare sotto rumore non serve a nulla di misurabile in questo
problema: verificato su due metriche, due punti di lavoro, e l'intera
griglia di rumore del Passo 5 --- non un'estrapolazione da un singolo
controllo.
\end{okbox}

\needspace{6\baselineskip}
\section*{File prodotti}
\addcontentsline{toc}{section}{File prodotti}

\texttt{vqe\_noise\_aware\_dimero.py} (ottimizzazione, controlli, secondo
punto di lavoro), \texttt{validate\_vqe\_noise\_aware\_dimero.py} (otto
controlli), \texttt{verifica\_passo5\_noise\_aware.py} (griglia
completa), questo documento. Derivazione quantitativa dell'argomento di
continuità in \texttt{nota\_continuita\_preparazione.tex} (nota di lavoro
interna). Aggiornamento in \texttt{log\_decisioni.md}.

\end{document}
